\documentclass[11pt]{article}
\usepackage[a4paper,margin=26mm]{geometry}
\usepackage{amsmath,amssymb,mathtools}
\usepackage{unicode-math}
\setmainfont{Latin Modern Roman}
\setmathfont{Latin Modern Math}
\usepackage{xurl,hyperref}
\hypersetup{hidelinks,pdftitle={Problem statement}}
\newcommand{\sref}[2]{\hyperlink{source:#1:#2}{[S#2]}}
\newcommand{\cataloguescope}[1]{#1}
\begin{document}
\section{Leaf spaces and transverse structures — Question 8.7}\label{statement}

\subsection{Definitions and mathematical statement}

\paragraph{Shared notation.}$\mathbb N=\{1,2,\ldots\}$, and $\mathbb Z,\mathbb Q,\mathbb R,\mathbb C$ denote the integers, rationals, reals and complex numbers. Any other conventions are specified locally.


A left order on a group $G$ is a total order $<$ satisfying $g<h\Rightarrow ag<ah$. Its positive cone is $G^+=\{g:1<g\}$, with $G=G^+\sqcup\{1\}\sqcup(G^+)^{-1}$ and $G^+G^+\subseteq G^+$. Fix finite word encodings for a finitely generated group. For a class $\mathsf T$ of machines, a $\mathsf T$-order means that its positive cone is recognized by a machine in that class; in the finite-state case the source considers geodesic word representatives.
\paragraph{Statement recorded in the supplied source.}
For fundamental groups $G=\pi_1(M)$ of hyperbolic $3$-manifolds, determine which computational kinds of left orders exist, as the machine class $\mathsf T$ varies. Include finite-state automata, Turing machines and Turing machines with a specified oracle $O$. Classify the complexity of recognizing positive elements, rather than only whether the group is left-orderable. \sref{statement}{2}
\subsection{Short English statement}
Classify the computational complexity of left orders on hyperbolic manifold groups, from finite automata to ordinary and oracle Turing machines.
\subsection{Sources}
\begin{description}
\item[\hypertarget{source:statement:1}{S1}] ULAM.AI and the UnsolvedMath Contributors, \emph{UnsolvedMath: A Curated Collection of Open Mathematics Problems}, catalogue entry. CC BY 4.0. \url{https://huggingface.co/datasets/ulamai/UnsolvedMath}.
\item[\hypertarget{source:statement:2}{S2}] Danny Calegari, \emph{Problems in foliations and laminations of 3-manifolds}, version 0.78 (2002), Question 8.7; Definitions 1.1--1.10 and local definitions and remarks. \url{https://arxiv.org/abs/math/0209081v1}.
\end{description}
\label{end:statement}
\end{document}
